\section{Conclusion}

In this paper, a lightweight and standardized framework for traffic light detection on resource-constrained edge devices using the AUTOSAR Adaptive Platform has been proposed. By integrating the YOLO11 architecture with a MobileNetV4 backbone and utilizing the LiteRT runtime, the system successfully addresses the computational bottlenecks inherent to general-purpose edge CPUs. This contribution directly supports the advancement of Edge AI for sensing and robotics by demonstrating the ability to deliver real-time perception without relying on specialized AI accelerators. The proposed solution achieves an inference latency of approximately 30--40 ms on a Raspberry Pi 5, representing a $\sim$2.4$\times$ speedup compared to the baseline YOLO11n. However, this efficiency entails a necessary trade-off in detection precision. While the lightweight model retains approximately 95 \% of the baseline accuracy for general detection (a relative reduction of merely 4.9\% in mAP@50), a notable limitation is observed in stricter localization metrics, with mAP@50-95 decreasing from 78.5\% to 72.3\%. Despite this limitation in bounding box tightness, the system maintains a stable throughput of 24--33 FPS, thereby validating the feasibility of deploying efficient Machine Learning for Sensor Systems within the cost-effective and scalable SDV ecosystem. Future research directions include investigating advanced lightweight models, such as hybrid Vision Transformers, and utilizing the Yocto Project for OS-level CPU optimizations to further reduce inference latency on both current and high-end edge devices.